% Regression test: \proofof and proofs nested in proofs, the shape a structured
% proof (pf2 and the like) writes, where a step of a proof is itself a proof.
% A nested proof used to take as its result the one that preceded the proof
% around it, so an enclosing \proofof did not reach it: the references of the
% inner proof went to the last result declared instead of to the pinned one.
% Three cases, all in one document:
%   * outer \proofof pinning both proofs (lem:a and lem:b -> thm:early), the
%     pin standing before the nested proof and the reference in it after;
%   * inner \proofof overriding the outer pin for the inner proof alone
%     (lem:c -> thm:early, lem:d -> thm:pinned);
%   * a claim stated inside a proof, whose own proof attaches to the claim and
%     not to the proof around it (lem:e -> clm:one), the case the inheritance
%     must not swallow; the claim itself supports the result being proved
%     (clm:one -> thm:claimed), being a step of its proof.
% thm:last is declared last on purpose and proves nothing: it is where the
% misattributed edges landed, so any edge reaching it is a regression.
% Uses only amsthm; single-pass and bibtex-free for determinism.
\documentclass{article}
\usepackage{amsthm}
\usepackage{proofgraph}

\newtheorem{theorem}{Theorem}
\newtheorem{lemma}[theorem]{Lemma}
\newtheorem{claim}[theorem]{Claim}

\begin{document}

\begin{lemma}\label{lem:a}A.\end{lemma}
\begin{lemma}\label{lem:b}B.\end{lemma}
\begin{lemma}\label{lem:c}C.\end{lemma}
\begin{lemma}\label{lem:d}D.\end{lemma}
\begin{lemma}\label{lem:e}E.\end{lemma}
\begin{theorem}\label{thm:early}Early.\end{theorem}
\begin{theorem}\label{thm:pinned}Pinned by the inner proof.\end{theorem}
\begin{theorem}\label{thm:last}Last declared, proves nothing.\end{theorem}

% The outer pin covers the nested proof too.
\begin{proof}
  \proofof{thm:early}
  By \ref{lem:a}.
  \begin{proof}
    And by \ref{lem:b}.
  \end{proof}
\end{proof}

% An inner pin wins for the inner proof, and for it alone.
\begin{proof}
  \proofof{thm:early}
  By \ref{lem:c}.
  \begin{proof}\proofof{thm:pinned}
    And by \ref{lem:d}.
  \end{proof}
\end{proof}

% A claim opened inside a proof takes its own proof back.
\begin{theorem}\label{thm:claimed}With a claim inside its proof.\end{theorem}
\begin{proof}
  \begin{claim}\label{clm:one}A claim.\end{claim}
  \begin{proof}
    By \ref{lem:e}.
  \end{proof}
\end{proof}

\end{document}
